\documentclass[10pt,a4paper]{amsart}
\usepackage{amsmath,amssymb,amsthm,graphicx}
\usepackage{verbatim}
\newenvironment{reference}{\comment}{\endcomment}
%\newenvironment{reference}{}{}
\newenvironment{slogan}{\comment}{\endcomment}
\newenvironment{history}{\comment}{\endcomment}

% For commutative diagrams we use Xy-pic
\usepackage[all]{xy}

% We use 2cell for 2-commutative diagrams.
\xyoption{2cell}
% Theorem environments.
%
\theoremstyle{plain}
\newtheorem{theorem}[subsection]{Theorem}
\newtheorem{proposition}[subsection]{Proposition}
\newtheorem{lemma}[subsection]{Lemma}

\theoremstyle{definition}
\newtheorem{definition}[subsection]{Definition}
\newtheorem{example}[subsection]{Example}
\newtheorem{exercise}[subsection]{Exercise}
\newtheorem{situation}[subsection]{Situation}

\theoremstyle{remark}
\newtheorem{remark}[subsection]{Remark}
\newtheorem{remarks}[subsection]{Remarks}

\numberwithin{equation}{subsection}

% Macros
%
\def\lim{\mathop{\mathrm{lim}}\nolimits}
\def\colim{\mathop{\mathrm{colim}}\nolimits}
\def\Spec{\mathop{\mathrm{Spec}}}
\def\Hom{\mathop{\mathrm{Hom}}\nolimits}
\def\Ext{\mathop{\mathrm{Ext}}\nolimits}
\def\SheafHom{\mathop{\mathcal{H}\!\mathit{om}}\nolimits}
\def\SheafExt{\mathop{\mathcal{E}\!\mathit{xt}}\nolimits}
\def\Sch{\mathit{Sch}}
\def\Mor{\mathop{\mathrm{Mor}}\nolimits}
\def\Ob{\mathop{\mathrm{Ob}}\nolimits}
\def\Sh{\mathop{\mathit{Sh}}\nolimits}
\def\NL{\mathop{N\!L}\nolimits}
\def\CH{\mathop{\mathrm{CH}}\nolimits}
\def\proetale{{pro\text{-}\acute{e}tale}}
\def\etale{{\acute{e}tale}}
\def\QCoh{\mathit{QCoh}}
\def\Ker{\mathop{\mathrm{Ker}}}
\def\Im{\mathop{\mathrm{Im}}}
\def\Coker{\mathop{\mathrm{Coker}}}
\def\Coim{\mathop{\mathrm{Coim}}}

% Boxtimes
%
\DeclareMathSymbol{\boxtimes}{\mathbin}{AMSa}{"02}

%
% Macros for moduli stacks/spaces
%
\def\QCohstack{\mathcal{QC}\!\mathit{oh}}
\def\Cohstack{\mathcal{C}\!\mathit{oh}}
\def\Spacesstack{\mathcal{S}\!\mathit{paces}}
\def\Quotfunctor{\mathrm{Quot}}
\def\Hilbfunctor{\mathrm{Hilb}}
\def\Curvesstack{\mathcal{C}\!\mathit{urves}}
\def\Polarizedstack{\mathcal{P}\!\mathit{olarized}}
\def\Complexesstack{\mathcal{C}\!\mathit{omplexes}}
% \Pic is the operator that assigns to X its picard group, usage \Pic(X)
% \Picardstack_{X/B} denotes the Picard stack of X over B
% \Picardfunctor_{X/B} denotes the Picard functor of X over B
\def\Pic{\mathop{\mathrm{Pic}}\nolimits}
\def\Picardstack{\mathcal{P}\!\mathit{ic}}
\def\Picardfunctor{\mathrm{Pic}}
\def\Deformationcategory{\mathcal{D}\!\mathit{ef}}
\usepackage{cleveref}
\crefname{originalreference}{original reference}{original references}
\Crefname{originalreference}{Original reference}{Original references}
\makeatletter
\newlabel{lemma-Rlim-quasi-coherent}{{0A0J}{source}{}{sourcepointer-1}{}}
\newlabel{lemma-Rlim-quasi-coherent@cref}{{[originalreference][0][]0A0J}{[1][0][]source}{}{}{}}
\makeatother
\crefname{originalreference}{original reference}{original references}
\Crefname{originalreference}{Original reference}{Original references}
\makeatletter
\newlabel{lemma-cohomology-base-change}{{08EU}{source}{}{sourcepointer-2}{}}
\newlabel{lemma-cohomology-base-change@cref}{{[originalreference][0][]08EU}{[1][0][]source}{}{}{}}
\makeatother
\crefname{originalreference}{original reference}{original references}
\Crefname{originalreference}{Original reference}{Original references}
\makeatletter
\newlabel{lemma-flat-proper-perfect-direct-image-general}{{0B91}{source}{}{sourcepointer-3}{}}
\newlabel{lemma-flat-proper-perfect-direct-image-general@cref}{{[originalreference][0][]0B91}{[1][0][]source}{}{}{}}
\makeatother
\crefname{originallemma}{original lemma}{original lemmas}
\Crefname{originallemma}{Original lemma}{Original lemmas}
\makeatletter
\newlabel{more-algebra-lemma-pseudo-coherent-tensor-limit}{{0D2L}{source}{}{sourcepointer-4}{}}
\newlabel{more-algebra-lemma-pseudo-coherent-tensor-limit@cref}{{[originallemma][0][]0D2L}{[1][0][]source}{}{}{}}
\makeatother
\crefname{originallemma}{original lemma}{original lemmas}
\Crefname{originallemma}{Original lemma}{Original lemmas}
\makeatletter
\newlabel{sites-cohomology-lemma-dual-perfect-complex}{{08JJ}{source}{}{sourcepointer-5}{}}
\newlabel{sites-cohomology-lemma-dual-perfect-complex@cref}{{[originallemma][0][]08JJ}{[1][0][]source}{}{}{}}
\makeatother
\crefname{originalsection}{original section}{original sections}
\Crefname{originalsection}{Original section}{Original sections}
\makeatletter
\newlabel{sites-cohomology-section-global-RHom}{{0B6E}{source}{}{sourcepointer-6}{}}
\newlabel{sites-cohomology-section-global-RHom@cref}{{[originalsection][0][]0B6E}{[1][0][]source}{}{}{}}
\makeatother
\crefname{originalreference}{original reference}{original references}
\Crefname{originalreference}{Original reference}{Original references}
\makeatletter
\newlabel{theorem-bondal-van-den-Bergh}{{09IS}{source}{}{sourcepointer-7}{}}
\newlabel{theorem-bondal-van-den-Bergh@cref}{{[originalreference][0][]09IS}{[1][0][]source}{}{}{}}
\makeatother
\title{Tensoring a proper flat finitely presented cover with an inverse limit of modules commutes up to DQ}
\author{The Stacks Project Authors}
\date{stacks-project:spaces-perfect.tex:local-snapshot-20261012; GNU FDL 1.2 or later}
\begin{document}
\maketitle
\noindent\textit{Editorial note.} All delivered mathematical text is the author's own, copied from the pinned Stacks Project chapter file spaces-perfect.tex; no statement or proof step was rewritten, reordered or removed.\par\medskip
\noindent\textit{Editorial note.} The counted proof is the author's complete argument: the proof reduces to perfect test objects via the perfect generator, converts the limit to a derived limit, moves RHom through it, applies the duality and base-change identifications, and closes with the pseudo-coherent limit commutation from More on Algebra\par\medskip
\noindent\textit{Editorial note.} The label lemma-pullback-and-limits names the counted statement, and the proof environment immediately following it in the source is delivered verbatim as its proof.\par\medskip
\noindent\textit{Editorial note.} The counted statement is carried by the author's own lemma environment in the chapter source.\par\medskip
\noindent\textit{Allowed standard prerequisites (editorial).}\begin{itemize}
\item the perfect generator of D\_QCoh(O\_X) for X proper, flat and of finite presentation over a ring (theorem-bondal-van-den-Bergh of the same chapter)
\item the identity lim O\_X tensor M\_n = Rlim O\_X tensor M\_n for systems with surjective transitions (lemma-Rlim-quasi-coherent of the same chapter)
\item RHom\_X(E, -) is exact and commutes with products, hence with derived limits (Cohomology on Sites, global RHom section)
\item the duality formula RHom\_X(E, O\_X tensor M) = RGamma(X, E\textasciicircum{}vee tensor\textasciicircum{}L Lf\textasciicircum{}*M) = RGamma(X, E\textasciicircum{}vee) tensor\textasciicircum{}L M for E perfect (sites-cohomology-lemma-dual-perfect-complex and lemma-cohomology-base-change of the same chapter)
\item RGamma(X, E\textasciicircum{}vee) is a perfect complex of R-modules for E perfect and f proper, flat, of finite presentation (lemma-flat-proper-perfect-direct-image-general of the same chapter)
\item pseudo-coherent complexes commute with tensoring and filtered limits (More on Algebra, Lemma more-algebra-lemma-pseudo-coherent-tensor-limit)
\end{itemize}\medskip
\ifdefined\hypertarget\hypertarget{sourcepointer-1}{}\fi
\noindent\textit{Source reference.} Original reference 0A0J: sources/stacks-spaces-perfect/pinned-original-source.tex, source lines 585--585. Prerequisite pointer invoked at proof line(s) 6512; it names a standard fact of the same chapter that the counted argument uses, and every step of the counted proof built on it is retained verbatim in the delivered body.\par\medskip
\ifdefined\hypertarget\hypertarget{sourcepointer-2}{}\fi
\noindent\textit{Source reference.} Original reference 08EU: sources/stacks-spaces-perfect/pinned-original-source.tex, source lines 4883--4883. Prerequisite pointer invoked at proof line(s) 6530; it names a standard fact of the same chapter that the counted argument uses, and every step of the counted proof built on it is retained verbatim in the delivered body.\par\medskip
\ifdefined\hypertarget\hypertarget{sourcepointer-3}{}\fi
\noindent\textit{Source reference.} Original reference 0B91: sources/stacks-spaces-perfect/pinned-original-source.tex, source lines 6442--6442. Prerequisite pointer invoked at proof line(s) 6531; it names a standard fact of the same chapter that the counted argument uses, and every step of the counted proof built on it is retained verbatim in the delivered body.\par\medskip
\ifdefined\hypertarget\hypertarget{sourcepointer-4}{}\fi
\noindent\textit{Source reference.} Original lemma 0D2L: sources/stacks-chow/tags.txt, source lines 16386--16386. Prerequisite pointer invoked at proof line(s) 6534; it names a standard fact of a sibling chapter that the counted argument uses, and every step of the counted proof built on it is retained verbatim in the delivered body.\par\medskip
\ifdefined\hypertarget\hypertarget{sourcepointer-5}{}\fi
\noindent\textit{Source reference.} Original lemma 08JJ: sources/stacks-chow/tags.txt, source lines 10852--10852. Prerequisite pointer invoked at proof line(s) 6523; it names a standard fact of a sibling chapter that the counted argument uses, and every step of the counted proof built on it is retained verbatim in the delivered body.\par\medskip
\ifdefined\hypertarget\hypertarget{sourcepointer-6}{}\fi
\noindent\textit{Source reference.} Original section 0B6E: sources/stacks-chow/tags.txt, source lines 14067--14067. Prerequisite pointer invoked at proof line(s) 6516; it names a standard fact of a sibling chapter that the counted argument uses, and every step of the counted proof built on it is retained verbatim in the delivered body.\par\medskip
\ifdefined\hypertarget\hypertarget{sourcepointer-7}{}\fi
\noindent\textit{Source reference.} Original reference 09IS: sources/stacks-spaces-perfect/pinned-original-source.tex, source lines 3891--3891. Prerequisite pointer invoked at proof line(s) 6510; it names a standard fact of the same chapter that the counted argument uses, and every step of the counted proof built on it is retained verbatim in the delivered body.\par\medskip
\begin{lemma}
\label{lemma-pullback-and-limits}
Let $R$ be a ring. Let $X$ be an algebraic space and let
$f : X \to \Spec(R)$ be proper, flat, and
of finite presentation. Let $(M_n)$ be an inverse
system of $R$-modules with surjective transition maps.
Then the canonical map
$$
\mathcal{O}_X \otimes_R (\lim M_n)
\longrightarrow
\lim \mathcal{O}_X \otimes_R M_n
$$
induces an isomorphism from the source to $DQ_X$ applied to the target.
\end{lemma}

\begin{proof}
The statement means that for any object $E$ of
$D_\QCoh(\mathcal{O}_X)$ the induced map
$$
\Hom(E, \mathcal{O}_X \otimes_R (\lim M_n))
\longrightarrow
\Hom(E, \lim \mathcal{O}_X \otimes_R M_n)
$$
is an isomorphism. Since $D_\QCoh(\mathcal{O}_X)$ has
a perfect generator (Theorem \ref{theorem-bondal-van-den-Bergh})
it suffices to check this for perfect $E$.
By Lemma \ref{lemma-Rlim-quasi-coherent} we have
$\lim \mathcal{O}_X \otimes_R M_n = R\lim \mathcal{O}_X \otimes_R M_n$.
The exact functor
$R\Hom_X(E, -) : D_\QCoh(\mathcal{O}_X) \to D(R)$
of Cohomology on Sites, Section \ref{sites-cohomology-section-global-RHom}
commutes with products and hence with derived limits, whence
$$
R\Hom_X(E, \lim \mathcal{O}_X \otimes_R M_n) =
R\lim R\Hom_X(E, \mathcal{O}_X \otimes_R M_n)
$$
Let $E^\vee$ be the dual perfect complex, see
Cohomology on Sites, Lemma \ref{sites-cohomology-lemma-dual-perfect-complex}.
We have
$$
R\Hom_X(E, \mathcal{O}_X \otimes_R M_n) =
R\Gamma(X, E^\vee \otimes_{\mathcal{O}_X}^\mathbf{L} Lf^*M_n) =
R\Gamma(X, E^\vee) \otimes_R^\mathbf{L} M_n
$$
by Lemma \ref{lemma-cohomology-base-change}.
From Lemma \ref{lemma-flat-proper-perfect-direct-image-general}
we see $R\Gamma(X, E^\vee)$ is a perfect complex of $R$-modules.
In particular it is a pseudo-coherent complex and by
More on Algebra, Lemma \ref{more-algebra-lemma-pseudo-coherent-tensor-limit}
we obtain
$$
R\lim R\Gamma(X, E^\vee) \otimes_R^\mathbf{L} M_n =
R\Gamma(X, E^\vee) \otimes_R^\mathbf{L} \lim M_n
$$
as desired.
\end{proof}

\end{document}
